\clearpage
\VerdierPhysicalPageAnchor{0061}%
\VerdierSemanticAnchor{ch1-1-6-2}
\VerdierSemanticAnchor{ch1-1-6-2-proposition}
\noindent\textbf{命題 1.6.2.}\quad

\VerdierSemanticAnchor{ch1-1-6-2-part1}
\noindent\emph{1) $\alpha\in Z^2(\mathbb Z^n,\mathbb Z/2\mathbb Z)$ を
正規化された非斉次余サイクル、
$\dot\alpha\in H^2(\mathbb Z^n,\mathbb Z/2\mathbb Z)$ をその
コホモロジー類、
$\dot\alpha_{i,j}\in H^2(\mathbb Z^n_{i,j},\mathbb Z/2\mathbb Z)$ を
その成分とする。このとき（(1.6.1.2) を用いて）：}

\VerdierSemanticAnchor{ch1-1-6-2-eq1}
\begin{equation}
  \dot\alpha_{i,j}
  =\alpha(e_i,e_j)^{-1}\alpha(e_j,e_i).
  \tag{1.6.2.1}
\end{equation}

\VerdierSemanticAnchor{ch1-1-6-2-convention}
\noindent\emph{（$\mathbb Z/2\mathbb Z$ の群演算は乗法的に記す。）}

\VerdierSemanticAnchor{ch1-1-6-2-part2}
\noindent\emph{2) $\alpha_1$ と $\alpha_2$ を
$Z^2(\mathbb Z^n,\mathbb Z/2\mathbb Z)$ のコホモロガスな二元とし、各
$i\in[1,n]$ に対して $\beta_i\in\mathbb Z/2\mathbb Z$ をとる。このとき、
写像}

\VerdierSemanticAnchor{ch1-1-6-2-beta-map}
\begin{equation*}
  \beta\colon\mathbb Z^n\longrightarrow\mathbb Z/2\mathbb Z
\end{equation*}
\noindent\emph{で次を満たすものがただ一つ存在する：}

\VerdierSemanticAnchor{ch1-1-6-2-eq2}
\begin{equation}
  \begin{aligned}
    \alpha_2(\sigma_1,\sigma_2)\alpha_1^{-1}(\sigma_1,\sigma_2)
      &=\beta(\sigma_2)\beta(\sigma_1+\sigma_2)^{-1}\beta(\sigma_1),\\[-1mm]
    &\hspace{7em}\text{ただし }\sigma_1,\sigma_2\in\mathbb Z^n,
  \end{aligned}
  \tag{1.6.2.2}
\end{equation}

\VerdierSemanticAnchor{ch1-1-6-2-eq3}
\begin{equation}
  \beta(e_i)=\beta_i.
  \tag{1.6.2.3}
\end{equation}

\VerdierSemanticAnchor{ch1-1-6-2-proof}
\VerdierSemanticAnchor{ch1-1-6-2-proof-first}
第1の主張を証明するため $n=2$ の場合に帰着し、$\alpha$ を余境界で変更しても
$\alpha(e_i,e_j)^{-1}\alpha(e_j,e_i)$ が変わらないことを直ちに確かめる。
$\alpha$ が余境界ならば対称であり、したがって公式 (1.6.2.1) はこの場合に
成り立つ。一方、余サイクル

\VerdierSemanticAnchor{ch1-1-6-2-alpha0}
\begin{equation*}
  \alpha_0(p_1e_i+q_1e_j,p_2e_i+q_2e_j)=(-1)^{q_1p_2}
\end{equation*}
は $0$ とコホモロガスではない。実際、

\VerdierSemanticAnchor{ch1-1-6-2-alpha0-test}
\begin{equation*}
  \alpha_0(e_i,e_j)^{-1}\alpha_0(e_j,e_i)=-1
\end{equation*}
であり、零とコホモロガスでない任意の余サイクルは $\alpha_0$ と
コホモロガスである。

\VerdierSemanticAnchor{ch1-1-6-2-proof-second}
第2の主張については、(1.6.2.2) の二つの解の比が群準同型であることに注意
する。したがって (1.6.2.2) の解は、各 $e_i$ 上で指定した値をとるように、
ただ一つの群準同型によって常に修正できる。

\VerdierSemanticAnchor{ch1-1-6-3}
\VerdierSemanticAnchor{ch1-1-6-3-definition}
\noindent\textbf{定義 1.6.3.}\quad
$e_1,\ldots,e_n$ を $\mathbb Z^n$ の基底とする。\emph{基底
$e_1,\ldots,e_n$ に関する Koszul 類}とは、各群
$H^2(\mathbb Z^n_{i,j},\mathbb Z/2\mathbb Z)$ への制限が非零である
$H^2(\mathbb Z^n,\mathbb Z/2\mathbb Z)$ のただ一つの元である。

\VerdierSemanticAnchor{ch1-1-6-3-basis-dependence}
Koszul 類の定義は $\mathbb Z^n$ の基底の選択に依存することに注意する。
一般には、異なる二つの基底に付随する Koszul 類は異なる。

\VerdierSemanticAnchor{ch1-1-6-3-filtration}
$e_1,\ldots,e_n$ を $[1,n]$ の自然順序に従って並べた $\mathbb Z^n$ の
基底とする。$j\leq i$ を満たす $e_j$ が生成する部分群を
$\mathbb Z^n(\leq i)$ と記し、$j\geq i$ を満たす $e_j$ が生成する部分群を
$\mathbb Z^n(\geq i)$ と記す。
